\section{Asynchronous vs Synchronous Execution Models}

\subsection{Ethereum: Synchronous EVM Calls}

Ethereum smart contracts execute within a synchronous, single-threaded environment.
When ContractA calls ContractB within the same transaction, the EVM pauses A's
execution, runs B to completion, and returns control to A. This model provides
\textbf{atomicity}: if any step in the chain reverts, the entire transaction
reverts, and all state changes are discarded.

\begin{itemize}
    \item \textbf{Deterministic ordering}: ContractB runs before ContractC,
          and the order is guaranteed by the EVM call stack.
    \item \textbf{Atomic rollback}: A failure in ContractC reverts B and A.
    \item \textbf{Reentrancy risk}: Synchronous calls create reentrancy vectors
          (e.g., the 2016 DAO hack), requiring careful use of checks-effects-interactions.
\end{itemize}

\subsection{Polkadot: Synchronous Cross-Contract Calls (ink! 5.0)}

ink! contracts on Polkadot also execute synchronously within a single dispatch.
When ContractA calls ContractB via \texttt{build\_call}, execution blocks until
B returns. The model closely mirrors Ethereum's in terms of atomicity:

\begin{itemize}
    \item \textbf{Same-block execution}: A $\to$ B $\to$ C all occur in one block.
    \item \textbf{Gas limits per call}: Each cross-contract call can specify its
          own gas limit, preventing one contract from exhausting another's budget.
    \item \textbf{Weight-based pricing}: Polkadot uses a weight system rather
          than gas, but the synchronous semantics are identical for the caller.
\end{itemize}

\subsection{TON: Asynchronous Message Passing}

TON fundamentally differs from Ethereum and Polkadot by adopting an
\textbf{asynchronous actor model}. Smart contracts (actors) communicate
exclusively via messages. When ContractA sends a message to ContractB:

\begin{itemize}
    \item ContractA does \textbf{not} wait for ContractB to process the message.
    \item The message is placed in a queue and may be processed in a different
          shard, by a different validator, in a different block.
    \item ContractA's execution completes immediately; there is no automatic
          rollback if ContractB fails.
    \item Error handling requires explicit \textbf{bounce messages}: if B fails,
          it can send a bounce message back to A, but A must handle it explicitly.
\end{itemize}

\subsection{Impact on the Chain-of-Contracts Pattern}

The choice of execution model directly affects how the A $\to$ B $\to$ C chain
must be designed:

\begin{table}[h]
\centering
\begin{tabular}{|l|c|c|c|}
\hline
\textbf{Property} & \textbf{Ethereum} & \textbf{Polkadot} & \textbf{TON} \\
\hline
Atomicity & Full & Full & None (by default) \\
Error propagation & Automatic revert & Automatic revert & Manual (bounce) \\
Call ordering & Guaranteed & Guaranteed & Not guaranteed \\
Shard-crossing & N/A & N/A & Required \\
\hline
\end{tabular}
\caption{Comparison of execution models for A$\to$B$\to$C chain}
\end{table}

For TON, implementing a reliable chain-of-contracts requires:
\begin{enumerate}
    \item Sending the message from A to B with a bounce flag.
    \item B sending its message to C (with B's remaining balance attached).
    \item C sending a confirmation message back through B to A.
    \item A implementing a timeout mechanism in case messages are lost.
\end{enumerate}

This increased complexity is the cost of TON's scalability advantages:
asynchronous messages enable infinite sharding without cross-shard coordination
overhead at execution time.